\centerline{\xiaoyi\bfseries Abstract}

Clinical decision-making is inherently longitudinal — clinicians must synthesize evidence across repeated hospitalizations, diagnostic tests, and evolving treatment responses spanning months or years. However, current benchmarks for medical large language model (LLM) agents predominantly evaluate single-encounter behavior, focusing on tool-use accuracy and fact retrieval rather than the cross-visit temporal integration that real clinical practice demands. This thesis presents the decision-making component of LongMedBench, a benchmark constructed from real-world Electronic Health Record (EHR) data (MIMIC-IV) that evaluates LLM agents on long-horizon clinical decision-making. We design a three-tier memory architecture operating at three granularities — Note Memory (coarse-grained visit summaries), Event Memory (fine-grained structured event sequences), and Contextual Memory (immediate visit context) — and systematically compare memory provision strategies including naive long-context, retrieval-augmented generation (RAG), and agent memory systems (Mem0). Through experiments on GPT-5-mini, DeepSeek-v3.2, and Qwen-Turbo, we find that decision-making accuracy across all models remains modest (approximately 0.55--0.62 exact match), and critically, that memory augmentation architectures provide no significant improvement over the naive long-context baseline. Decision accuracy depends primarily on the immediate clinical context, with historical information contributing negligible marginal value regardless of memory granularity or retrieval strategy. This central finding reveals a retrieval-reasoning gap: current architectures can locate relevant historical facts but cannot effectively integrate them into proactive clinical planning. We discuss the implications of this gap for medical AI agent design and propose directions for memory architectures that go beyond retrieval to support structured longitudinal reasoning.

\vspace{2em}
\noindent\textbf{Keywords:} Medical AI Agents, Clinical Decision-Making, EHR Benchmark, Long-Horizon Reasoning, Memory-Augmented LLM, Retrieval-Augmented Generation
